\chapter{Introduzione}
\label{cap:introduzione}

\section{Il problema: dal testo al gloss}

La lingua dei segni americana (\emph{American Sign Language}, ASL) è una lingua
naturale completa, con una propria grammatica, che non coincide con l'inglese.
Non esiste una forma scritta ufficiale dell'ASL. Per lavorare con strumenti
testuali si usa quindi una convenzione di trascrizione chiamata \emph{gloss}:
ogni segno viene rappresentato dalla parola inglese in maiuscolo che gli
corrisponde più da vicino, eventualmente con marcatori aggiuntivi.

Un esempio reale tratto dal corpus impiegato in questo lavoro:

\begin{center}
\begin{tabular}{ll}
\toprule
Inglese & \texttt{they will not flinch .} \\
Gloss   & \texttt{X-Y WILL DESC-NOT FLINCH .} \\
\midrule
Inglese & \texttt{this kind of thing must never happen again .} \\
Gloss   & \texttt{THIS KIND THING MUST DESC-NEVER HAPPEN DESC-AGAIN .} \\
\bottomrule
\end{tabular}
\end{center}

Va detto subito, perché condiziona tutta l'interpretazione dei risultati: il
gloss \textbf{non è} la lingua dei segni. È una trascrizione lessicale
impoverita, che perde la struttura spaziale, i marcatori non manuali (espressioni
facciali, direzione dello sguardo) e la simultaneità tipica del segnato. Parte
della letteratura sostiene che il gloss sia una rappresentazione
inadeguata: Yin e Read~\cite{yinread2020stmc} osservano che nei loro esperimenti
la traduzione a partire dal \emph{video} batte quella a partire dal gloss di
riferimento, e concludono che il gloss è una rappresentazione inefficiente della
lingua dei segni. Sul piano più generale del trattamento delle lingue dei segni
nell'elaborazione del linguaggio naturale, Yin et
al.~\cite{yin2021including} raccolgono le raccomandazioni metodologiche del
settore, fra cui l'avvertenza a non trattare il gloss come se fosse la lingua.

Nonostante questi limiti il gloss resta l'interfaccia testuale su cui si
costruiscono i sistemi esistenti: esistono baseline aperte basate su gloss per
la traduzione da lingua parlata a lingua dei segni~\cite{moryossef2023baseline} e
architetture di produzione che trattano il compito come selezione e riordino di
segni~\cite{walsh2024select}.

Il compito affrontato qui, detto \emph{text-to-gloss} (T2G), è dunque:

\begin{equation}
  \text{inglese scritto} \;\longmapsto\; \text{sequenza di gloss ASL}
\end{equation}

Si tratta di un passo intermedio in una ipotetica pipeline più ampia
(testo $\to$ gloss $\to$ avatar segnante), non di un sistema di traduzione
completo.

\section{Che cosa è stato costruito}

Il sistema si basa su un modello linguistico di piccole dimensioni,
\texttt{Qwen2.5-0.5B-Instruct}, adattato al compito con due famiglie di metodi:

\begin{itemize}
  \item \textbf{SFT} (\emph{supervised fine-tuning}): addestramento supervisionato
        classico, in cui il modello impara a riprodurre il gloss di riferimento;
  \item \textbf{GRPO} (\emph{Group Relative Policy Optimization}): apprendimento
        per rinforzo, in cui il modello genera più candidati per ogni frase e
        viene premiato in base alla qualità relativa di ciascuno.
\end{itemize}

Le due famiglie sono state combinate in quattro celle sperimentali: modello base
senza addestramento, solo SFT, solo GRPO a partire dal modello base, e la
pipeline completa SFT seguito da GRPO.

A queste si sovrappongono due fattori ortogonali:

\begin{itemize}
  \item la modalità di \emph{prompting}: \emph{zero-shot} (solo la frase da
        tradurre) oppure \emph{few-shot} con retrieval (nel prompt vengono
        inseriti $k=3$ esempi simili recuperati dal training set);
  \item la \emph{decodifica vincolata}: un componente simbolico (un
        \emph{Trie}, ossia un albero di prefissi) che a ogni passo di generazione
        vieta al modello di emettere token che non portino a un gloss valido.
\end{itemize}

Quest'ultimo elemento è ciò che rende il sistema \emph{neuro-simbolico}: una rete
neurale produce le probabilità, una struttura dati simbolica ne restringe il
supporto.

\section{Il risultato principale, in anticipo}

Questa relazione non riporta un miglioramento delle metriche. Riporta un
\textbf{risultato negativo con meccanismo}, che è più utile di un miglioramento
apparente perché spiega \emph{perché} la domanda iniziale era mal posta.

Il fatto centrale, misurato e documentato nel
Capitolo~\ref{cap:dataset}, è che il corpus impiegato (ASLG-PC12) non è stato
prodotto da annotatori umani: è \textbf{sintetico}, generato applicando regole
all'inglese. La conseguenza pratica è che una baseline a regole, costruita in
circa quindici righe di codice contando le corrispondenze parola-per-parola nel
solo training set, ottiene:

\begin{center}
\begin{tabular}{lrr}
\toprule
Sistema & ROUGE-L & Exact match \\
\midrule
Baseline a regole (nessun apprendimento) & $0{,}9697$ & $0{,}5912$ \\
SFT (modello addestrato) & $0{,}9752$ & $0{,}8130$ \\
GRPO da modello base & $0{,}6076$ & $0{,}0528$ \\
\bottomrule
\end{tabular}
\end{center}

Si legga con attenzione la terza riga: il miglior risultato ottenuto con
reinforcement learning ($0{,}6076$) sta \emph{molto sotto} una regola che non
apprende nulla ($0{,}9697$). Persino la trasformazione banale ``metti tutto in
maiuscolo'' raggiunge ROUGE-L $0{,}6454$, cioè batte il GRPO.

Da qui seguono tre conseguenze che strutturano il resto della relazione:

\begin{enumerate}
  \item \textbf{Nessun numero è interpretabile senza la baseline a regole
        accanto.} Un risultato di $0{,}6$ su questo corpus non è un successo
        parziale: è un fallimento rispetto a un riferimento che nessuno aveva
        calcolato.
  \item \textbf{La domanda ``il GRPO aiuta il T2G?'' non è ben posta su
        ASLG-PC12.} Il compito è quasi deterministico, l'SFT lo satura, e
        qualsiasi metodo di rinforzo parte già in svantaggio rispetto a una
        regola. La domanda si riduce a ``il GRPO memorizza una mappa
        quasi-identità meglio dell'SFT?'', che riguarda l'ottimizzazione e non
        la lingua dei segni.
  \item \textbf{La metrica di riferimento è essa stessa difettosa.} Come
        mostrato nel Capitolo~\ref{cap:metriche}, l'implementazione di ROUGE-L
        impiegata è insensibile alle maiuscole e spezza i token sui caratteri
        non alfanumerici, quindi premia un modello che si limita a copiare
        l'inglese.
\end{enumerate}

\section{Contributi}

Al netto del risultato negativo, il lavoro produce alcuni contributi
verificabili:

\begin{itemize}
  \item una \textbf{baseline a regole} riutilizzabile
        (\texttt{src/analysis/rule\_baseline.py}) che rende interpretabili i
        numeri su questo corpus, insieme a una metrica insensibile alla copia
        (\emph{non-copy-token accuracy});
  \item un \textbf{censimento degli artefatti} del generatore del corpus: il
        $10{,}47\%$ delle righe di training contiene parole corrotte da una
        rimozione di sottostringa (per esempio \texttt{therefore} diventa
        \texttt{DESC-REFORE});
  \item la \textbf{quantificazione del supporto dei rollout}, che spiega
        meccanicamente quando il reinforcement learning non può funzionare
        (Capitolo~\ref{cap:grpo});
  \item un risultato \textbf{neuro-simbolico non ovvio}: il vincolo simbolico
        peggiora la metrica di sovrapposizione e contemporaneamente \emph{abilita}
        il segnale di reward (Capitolo~\ref{cap:decoding});
  \item due \textbf{obiettivi ausiliari} implementati e qualificati prima di
        essere addestrati, con i relativi criteri di promozione
        (Capitolo~\ref{cap:ausiliari}).
\end{itemize}

\section{Come leggere questa relazione}

La relazione è pensata per essere autocontenuta. Il
Capitolo~\ref{cap:fondamenti} introduce da zero le nozioni teoriche necessarie:
modelli linguistici autoregressivi, entropia incrociata, adattamento a rango
basso, apprendimento per rinforzo, gradiente di policy. Chi ha già familiarità
con questi argomenti può saltarlo.

I capitoli successivi seguono la pipeline: dati
(Capitolo~\ref{cap:dataset}), modello (Capitolo~\ref{cap:modello}), i due
metodi di addestramento (Capitoli~\ref{cap:sft} e~\ref{cap:grpo}), il disegno
delle reward (Capitolo~\ref{cap:reward}), il componente simbolico
(Capitolo~\ref{cap:decoding}), gli obiettivi ausiliari
(Capitolo~\ref{cap:ausiliari}), le metriche (Capitolo~\ref{cap:metriche}) e i
risultati (Capitolo~\ref{cap:risultati}). Gli ultimi tre capitoli riguardano
l'infrastruttura di calcolo, l'organizzazione del codice e le conclusioni.

Ogni capitolo segue lo stesso schema: prima la teoria, poi la scelta
implementativa, poi il file e la funzione che la realizzano. I valori numerici
riportati sono tutti misurati; dove un valore non è disponibile è indicato
esplicitamente con un segnaposto, invece di essere stimato.
